% Partie entière et régularité : recollements de fonctions régulières
% Compile avec pdflatex, xelatex ou lualatex.
\documentclass[11pt,a4paper]{article}

\usepackage{iftex}
\ifPDFTeX
  \usepackage[utf8]{inputenc}
  \usepackage[T1]{fontenc}
  \IfFileExists{lmodern.sty}{\usepackage{lmodern}}{}
  \usepackage[french]{babel}
\else
  \usepackage{fontspec}
  \usepackage{polyglossia}
  \setdefaultlanguage{french}
\fi
\usepackage[margin=2.6cm]{geometry}
\usepackage{amsmath,amssymb,amsthm,mathtools}
\usepackage{enumitem}
\usepackage{microtype}
\usepackage[hidelinks]{hyperref}

% ---------- Environnements ----------
\theoremstyle{plain}
\newtheorem{theoreme}{Théorème}[section]
\newtheorem{proposition}[theoreme]{Proposition}
\newtheorem{lemme}[theoreme]{Lemme}
\newtheorem{corollaire}[theoreme]{Corollaire}
\theoremstyle{definition}
\newtheorem{definition}[theoreme]{Définition}
\newtheorem{exemple}[theoreme]{Exemple}
\newtheorem{exercice}{Exercice}
\theoremstyle{remark}
\newtheorem{remarque}[theoreme]{Remarque}

% ---------- Notations ----------
\newcommand{\N}{\mathbb{N}}
\newcommand{\Z}{\mathbb{Z}}
\newcommand{\R}{\mathbb{R}}
\newcommand{\Cc}{\mathcal{C}}
\newcommand{\1}{\mathbf{1}}
\DeclareMathOperator{\ord}{ord}

\title{Partie entière et régularité\\[2pt]
\large Pourquoi $E(x)\sin(\pi x)$ est continue, et comment la rendre plus lisse}
\author{Michel Hua}
\date{Octobre 2026}

\begin{document}
\maketitle

\begin{abstract}
En terminale, on étudie $f(x)=E(x)\sin(\pi x)$ sur $[-1,2]$. On trouve qu'elle est continue, alors que la partie entière $E$ saute à chaque entier, et qu'elle n'est pas dérivable en $0$ ni en $1$. Le calcul est facile, mais il laisse une question : \emph{pourquoi} le facteur $\sin(\pi x)$ efface-t-il les sauts, et pourquoi n'efface-t-il pas les points anguleux ? On répond avec une seule idée : un saut multiplié par une fonction qui s'annule à l'ordre $m$ donne une fonction de classe $\Cc^{m-1}$, et pas mieux. On en déduit que $E(x)\sin^m(\pi x)$ est exactement de classe $\Cc^{m-1}$. On termine par le théorème de Darboux, qui explique pourquoi une dérivée ne peut jamais sauter.
\end{abstract}

\paragraph{Leçons concernées.} 228 (Continuité, dérivabilité des fonctions réelles d'une variable réelle), 229 (Fonctions monotones, fonctions convexes), 218 (Formules de Taylor).

\paragraph{Développements possibles.} Lemme de troncature \ref{lem:troncature} et théorème de recollement \ref{thm:recollement}, avec le corollaire \ref{cor:fm} ; théorème de Darboux \ref{thm:darboux} avec la proposition \ref{prop:limite-derivee} et l'exemple \ref{ex:x2sin}. Une formalisation Lean de la fiche (\texttt{PartieEntiereRegularite.lean}) est en cours.

\paragraph{Prérequis.} Limites à gauche et à droite, dérivée, théorème des accroissements finis, formule de Leibniz pour la dérivée $n$-ième d'un produit.

\paragraph{Notations.} $E(x)$ désigne la partie entière de $x\in\R$, l'unique entier tel que $E(x)\leq x<E(x)+1$. Pour une fonction $g$ et un point $a$, on note $g(a^-)$ et $g(a^+)$ les limites de $g(x)$ quand $x\to a$ avec $x<a$, respectivement $x>a$. On note $\varphi'_g(a)$ et $\varphi'_d(a)$ les dérivées à gauche et à droite de $\varphi$ en $a$. Pour $a\in\R$, $\1_{]-\infty,a[}$ est l'indicatrice de $]-\infty,a[$. Par convention, « de classe $\Cc^{-1}$ » ne demande rien.

\tableofcontents

% ==========================================================
\section{La question de départ}
% ==========================================================

On commence par faire l'exercice de terminale, à la main. Cela fixe les valeurs qu'il faudra expliquer ensuite.

La partie entière est constante sur chaque intervalle $[k,k+1[$, $k\in\Z$, et elle saute de $1$ en chaque entier :
\[
  E(k^-)=k-1, \qquad E(k^+)=E(k)=k .
\]
Le facteur $\sin(\pi x)$, lui, s'annule exactement aux entiers. On s'attend donc à ce que les deux phénomènes se compensent.

\begin{exemple}\label{ex:origine}
Soit $f(x)=E(x)\sin(\pi x)$ pour $x\in[-1,2]$. On écrit $f$ morceau par morceau :
\[
  f(x)=\begin{cases}
    -\sin(\pi x) & \text{si } -1\leq x<0,\\
    0 & \text{si } 0\leq x<1,\\
    \sin(\pi x) & \text{si } 1\leq x<2,\\
    2\sin(2\pi)=0 & \text{si } x=2.
  \end{cases}
\]
\begin{itemize}[nosep]
\item \emph{Continuité.} Sur chacun des intervalles $[-1,0[$, $[0,1[$, $[1,2[$, $f$ est la restriction d'une fonction continue. Donc $f$ est continue en tout point de $]-1,0[\cup]0,1[\cup]1,2[$, et continue à droite en $-1$. Il reste les points de raccord. En $0$ : $f(0^-)=-\sin0=0=f(0)$. En $1$ : $f(1^-)=0=\sin\pi=f(1)$. En $2$ : $f(2^-)=\sin(2\pi)=0=f(2)$. Donc $f$ est continue sur $[-1,2]$.
\item \emph{Dérivabilité en $0$.} À gauche, $f=-\sin(\pi\,\cdot)$, donc $f'_g(0)=-\pi\cos0=-\pi$. À droite, $f$ est nulle, donc $f'_d(0)=0$. Comme $-\pi\neq0$, $f$ n'est pas dérivable en $0$. La courbe a un point anguleux : une demi-tangente de pente $-\pi$ à gauche, horizontale à droite.
\item \emph{Dérivabilité en $1$.} À gauche, $f$ est nulle, donc $f'_g(1)=0$. À droite, $f=\sin(\pi\,\cdot)$, donc $f'_d(1)=\pi\cos\pi=-\pi$. Là encore, $f$ n'est pas dérivable en $1$.
\end{itemize}
\emph{Variations sur $[1,2]$.} Comme $f(2)=0=\sin(2\pi)$, on a $f(x)=\sin(\pi x)$ sur tout le segment $[1,2]$. Donc $f$ y est dérivable (à droite en $1$, à gauche en $2$), avec $f'(x)=\pi\cos(\pi x)$.
\begin{itemize}[nosep]
\item Pour $x\in[1,\tfrac32]$, $\pi x\in[\pi,\tfrac{3\pi}2]$, donc $\cos(\pi x)\leq0$, avec égalité seulement en $x=\tfrac32$.
\item Pour $x\in[\tfrac32,2]$, $\pi x\in[\tfrac{3\pi}2,2\pi]$, donc $\cos(\pi x)\geq0$, avec égalité seulement en $x=\tfrac32$.
\end{itemize}
Ainsi $f$ est strictement décroissante sur $[1,\tfrac32]$, de $f(1)=0$ à $f(\tfrac32)=\sin\tfrac{3\pi}2=-1$, puis strictement croissante sur $[\tfrac32,2]$, de $-1$ à $f(2)=0$. Son minimum sur $[1,2]$ est $-1$, atteint seulement en $\tfrac32$.
\end{exemple}

\begin{remarque}
Sur $[-1,0]$, $f(x)=-\sin(\pi x)\geq0$, avec un maximum $f(-\tfrac12)=1$. Sur $[0,1]$, $f$ est nulle. Le graphe de $f$ sur $[-1,2]$ est donc formé de deux arches de sinusoïde, l'une au-dessus de l'axe, l'autre au-dessous, séparées par un palier nul. Le facteur $E(x)$ règle la hauteur des arches : sur $[k,k+1]$, l'arche a pour hauteur $|k|$.
\end{remarque}

Le calcul est terminé, mais il ne dit pas \emph{pourquoi} les choses se passent ainsi. On se pose trois questions.
\begin{enumerate}[nosep]
\item Pourquoi le produit est-il continu, alors que $E$ ne l'est pas ?
\item D'où viennent les pentes $-\pi$ et $0$ ? Peut-on les prévoir sans calcul ?
\item Peut-on, en changeant le facteur $\sin(\pi x)$, obtenir une fonction dérivable, ou même de classe $\Cc^1$, $\Cc^2$, etc. ?
\end{enumerate}
Les sections 2 et 3 répondent à ces questions. La section 4 explique pourquoi, dans la question 3, la régularité s'arrête toujours d'un cran trop tôt.

% ==========================================================
\section{Un saut multiplié par une fonction qui s'annule}
% ==========================================================

On isole ce qui se passe en \emph{un} point de saut $a$. La fonction à sauts $g$ joue le rôle de $E$, et la fonction $h$ celui de $\sin(\pi x)$. On ne demande presque rien à $g$ : seulement d'avoir une limite à gauche et une limite à droite.

L'idée est simple. Près de $a$, $g$ reste bornée et $h$ tend vers $0$, donc le produit tend vers $0$. Pour la dérivée, on écrit le taux d'accroissement du produit comme $g(x)$ multiplié par le taux d'accroissement de $h$ : le premier facteur tend vers $g(a^-)$ ou $g(a^+)$ selon le côté, le second vers $h'(a)$.

\begin{proposition}\label{prop:produit}
Soit $I$ un intervalle ouvert, $a\in I$, et $g,h\colon I\to\R$. On suppose que $g$ admet des limites finies $g(a^-)$ et $g(a^+)$, et que $h(a)=0$. On pose $\varphi=gh$.
\begin{enumerate}[nosep]
\item Si $h$ est continue en $a$, alors $\varphi$ est continue en $a$.
\item Si $h$ est dérivable en $a$, alors $\varphi$ est dérivable à gauche et à droite en $a$, avec
\[
  \varphi'_g(a)=g(a^-)\,h'(a), \qquad \varphi'_d(a)=g(a^+)\,h'(a).
\]
\item Dans ce cas, $\varphi$ est dérivable en $a$ si et seulement si $h'(a)=0$ ou $g(a^-)=g(a^+)$.
\end{enumerate}
\end{proposition}

\begin{proof}
On a $\varphi(a)=g(a)h(a)=0$, quelle que soit la valeur $g(a)$. Comme $g$ a des limites finies à gauche et à droite en $a$, elle est bornée sur $]a-\eta,a+\eta[\setminus\{a\}$ pour un $\eta>0$, disons par $M$.

1. Pour $0<|x-a|<\eta$, on a $|\varphi(x)-\varphi(a)|=|g(x)|\,|h(x)|\leq M|h(x)|$, qui tend vers $0$.

2. Pour $x\neq a$, comme $h(a)=0$,
\[
  \frac{\varphi(x)-\varphi(a)}{x-a}=g(x)\,\frac{h(x)-h(a)}{x-a}.
\]
Ce produit tend vers $g(a^-)h'(a)$ quand $x\to a^-$, et vers $g(a^+)h'(a)$ quand $x\to a^+$.

3. Une fonction est dérivable en $a$ si et seulement si ses dérivées à gauche et à droite existent et sont égales. Ici, cela s'écrit $\big(g(a^+)-g(a^-)\big)h'(a)=0$.
\end{proof}

\begin{remarque}
La formule se lit facilement : le défaut de dérivabilité
\[
  \varphi'_d(a)-\varphi'_g(a)=\big(g(a^+)-g(a^-)\big)\,h'(a)
\]
est le \emph{saut de $g$} multiplié par la \emph{pente de $h$}. Pour que le point anguleux disparaisse, il faut que l'un des deux soit nul. L'hypothèse $h(a)=0$, elle, est indispensable dès que $g$ saute : si $h$ est continue en $a$ avec $h(a)\neq0$, alors $\varphi(a^-)=g(a^-)h(a)$ et $\varphi(a^+)=g(a^+)h(a)$ diffèrent, et $\varphi$ n'est pas continue en $a$.
\end{remarque}

On applique maintenant la proposition à $E$ et à $\sin(\pi x)$, en tous les entiers à la fois.

\begin{corollaire}\label{cor:Esin}
La fonction $F\colon x\mapsto E(x)\sin(\pi x)$ est continue sur $\R$ et de classe $\Cc^\infty$ sur $\R\setminus\Z$. En tout $k\in\Z$,
\[
  F'_g(k)=(-1)^k(k-1)\pi, \qquad F'_d(k)=(-1)^k k\pi .
\]
En particulier, $F$ n'est dérivable en aucun entier.
\end{corollaire}

\begin{proof}
Sur $]k,k+1[$, $F$ coïncide avec $x\mapsto k\sin(\pi x)$, d'où la régularité hors de $\Z$. En $k\in\Z$, on applique la proposition \ref{prop:produit} avec $g=E$ et $h=\sin(\pi\,\cdot)$. On a $h(k)=0$, $h'(k)=\pi\cos(\pi k)=(-1)^k\pi$, $E(k^-)=k-1$ et $E(k^+)=k$. Les deux dérivées diffèrent de $(-1)^k\pi\neq0$.
\end{proof}

Pour $k=0$, on retrouve $F'_g(0)=-\pi$ et $F'_d(0)=0$ ; pour $k=1$, $F'_g(1)=0$ et $F'_d(1)=-\pi$. Ce sont les valeurs de l'exemple \ref{ex:origine}. Les deux premières questions ont leur réponse :
\begin{itemize}[nosep]
\item le produit est continu parce que $\sin(\pi x)$ s'annule là où $E$ saute ;
\item il y a un point anguleux parce que $\sin(\pi x)$ s'annule avec une pente non nulle, $\pm\pi$, et l'écart entre les deux pentes est le saut de $E$, qui vaut $1$, multiplié par cette pente.
\end{itemize}

% ==========================================================
\section{Régularité exacte d'un recollement}
% ==========================================================

On passe à la troisième question. Pour supprimer le point anguleux, la remarque de la section 2 dit qu'il faut un facteur qui s'annule \emph{avec une pente nulle}. Le plus simple est d'élever au carré : $\sin^2(\pi x)$ s'annule aux entiers avec une dérivée nulle. Plus généralement, $\sin^m(\pi x)$ s'annule « à l'ordre $m$ ». On va montrer que chaque ordre d'annulation supplémentaire fait gagner exactement un cran de régularité.

Pour mesurer « à quel ordre » une fonction s'annule, on compte combien de ses dérivées sont nulles.

\begin{definition}
Soit $s$ de classe $\Cc^\infty$ au voisinage de $a$. L'\emph{ordre d'annulation} de $s$ en $a$ est
\[
  \ord_a(s)=\min\{j\in\N : s^{(j)}(a)\neq0\}\in\N\cup\{+\infty\},
\]
avec $\ord_a(s)=+\infty$ si toutes les dérivées de $s$ sont nulles en $a$.
\end{definition}

Le cas modèle est celui d'une fonction qu'on « coupe » en $a$ : on garde $s$ à gauche et on met $0$ à droite. C'est un saut de hauteur $1$ (l'indicatrice) multiplié par $s$. La proposition \ref{prop:produit} traite une dérivée ; on l'applique autant de fois qu'il y a de dérivées nulles.

\begin{lemme}[Troncature]\label{lem:troncature}
Soit $I$ un intervalle ouvert, $a\in I$, $m\geq1$, et $s\in\Cc^m(I)$ telle que $s(a)=s'(a)=\cdots=s^{(m-1)}(a)=0$. On pose $u=s\,\1_{]-\infty,a[}$. Alors $u$ est de classe $\Cc^{m-1}$ sur $I$, avec $u^{(j)}=s^{(j)}\1_{]-\infty,a[}$ pour $0\leq j\leq m-1$, et
\[
  \big(u^{(m-1)}\big)'_g(a)=s^{(m)}(a), \qquad \big(u^{(m-1)}\big)'_d(a)=0 .
\]
\end{lemme}

\begin{proof}[Idée de la preuve]
À chaque étape, $u^{(j)}$ est encore de la forme « indicatrice fois une fonction nulle en $a$ », et la proposition \ref{prop:produit} s'applique. Tant que la dérivée suivante de $s$ est nulle en $a$, les deux pentes sont égales et l'on peut continuer. À la dernière étape, la pente $s^{(m)}(a)$ n'est plus nulle et l'on s'arrête.
\end{proof}

\begin{proof}
Notons $\chi=\1_{]-\infty,a[}$, qui vérifie $\chi(a^-)=1$ et $\chi(a^+)=0$. Hors de $a$, $u$ coïncide localement avec $s$ ou avec $0$, donc tout se joue en $a$. On montre par récurrence sur $j\in\{0,\dots,m-1\}$ que $u$ est $j$ fois dérivable sur $I$, avec $u^{(j)}=s^{(j)}\chi$ continue.

Pour $j=0$ : $u=s\chi$ est continue en $a$ par la proposition \ref{prop:produit}(1), appliquée à $g=\chi$ et $h=s$, puisque $s(a)=0$.

Supposons le résultat vrai au rang $j\leq m-2$. La proposition \ref{prop:produit}(2), avec $g=\chi$ et $h=s^{(j)}$ (dérivable, nulle en $a$), montre que $u^{(j)}=s^{(j)}\chi$ admet en $a$ la dérivée à gauche $s^{(j+1)}(a)$ et la dérivée à droite $0$. Comme $j+1\leq m-1$, on a $s^{(j+1)}(a)=0$. Donc $u^{(j)}$ est dérivable en $a$, de dérivée $0=s^{(j+1)}(a)\chi(a)$. Hors de $a$, la dérivée de $s^{(j)}\chi$ est $s^{(j+1)}\chi$. Ainsi $u^{(j+1)}=s^{(j+1)}\chi$, qui est continue en $a$ par la proposition \ref{prop:produit}(1), puisque $s^{(j+1)}(a)=0$.

Enfin, au rang $m-1$, la fonction $h=s^{(m-1)}$ est dérivable en $a$ car $s\in\Cc^m(I)$, et nulle en $a$. La proposition \ref{prop:produit}(2) donne pour $u^{(m-1)}=s^{(m-1)}\chi$ la dérivée à gauche $s^{(m)}(a)\cdot1$ et la dérivée à droite $s^{(m)}(a)\cdot0=0$.
\end{proof}

Un recollement quelconque de deux fonctions lisses se ramène à ce cas modèle : on écrit la fonction recollée comme la fonction de droite, plus la différence des deux fonctions coupée en $a$.

\begin{theoreme}[Recollement]\label{thm:recollement}
Soit $I$ un intervalle ouvert, $a\in I$, $p,q\in\Cc^\infty(I)$, et
\[
  \varphi(x)=\begin{cases} p(x) & \text{si } x<a,\\ q(x) & \text{si } x\geq a.\end{cases}
\]
Soit $r=\ord_a(p-q)$. Alors $\varphi$ est de classe $\Cc^{r-1}$ sur $I$. Si $r$ est fini, $\varphi^{(r-1)}$ admet en $a$ la dérivée à gauche $p^{(r)}(a)$ et la dérivée à droite $q^{(r)}(a)$, qui sont distinctes ; en particulier, $\varphi$ n'est pas $r$ fois dérivable en $a$. Autrement dit, pour $n\in\N$,
\[
  \varphi\in\Cc^n(I)\iff (p-q)^{(j)}(a)=0 \text{ pour } 0\leq j\leq n .
\]
\end{theoreme}

\begin{proof}
Posons $s=p-q$ et $\chi=\1_{]-\infty,a[}$. Alors $\varphi=q+s\chi$, et $q$ est de classe $\Cc^\infty$ : toute la difficulté est dans $s\chi$.

Si $r=0$, il n'y a rien à montrer pour la classe $\Cc^{-1}$, et $\varphi$ n'est pas continue en $a$ puisque $\varphi(a^-)=p(a)\neq q(a)=\varphi(a)$.

Si $1\leq r<\infty$, le lemme \ref{lem:troncature} avec $m=r$ montre que $s\chi$ est de classe $\Cc^{r-1}$, donc aussi $\varphi$. Il montre aussi que $\varphi^{(r-1)}=q^{(r-1)}+(s\chi)^{(r-1)}$ admet en $a$ la dérivée à gauche $q^{(r)}(a)+s^{(r)}(a)=p^{(r)}(a)$ et la dérivée à droite $q^{(r)}(a)$. Elles diffèrent de $s^{(r)}(a)\neq0$.

Si $r=+\infty$, le lemme s'applique pour tout $m\geq1$, donc $\varphi$ est de classe $\Cc^\infty$.

Pour l'équivalence : la condition de droite signifie $r\geq n+1$. Si elle est vérifiée, $\varphi$ est de classe $\Cc^{r-1}$, donc $\Cc^n$. Sinon, $r\leq n$ et $\varphi$ n'est pas $r$ fois dérivable en $a$, donc pas de classe $\Cc^n$.
\end{proof}

\begin{remarque}
Le théorème dit que la régularité d'un recollement se lit sur \emph{une seule} fonction, la différence $p-q$, et en \emph{un seul} point. Il n'y a pas d'intermédiaire : si les $n+1$ premières dérivées se raccordent, la fonction est $\Cc^n$ ; dès qu'une dérivée ne se raccorde pas, tout s'arrête à ce cran.
\end{remarque}

Pour l'appliquer à $E(x)\sin^m(\pi x)$, il reste à vérifier que $\sin^m(\pi x)$ s'annule à l'ordre exactement $m$ en chaque entier. L'idée : près de $k$, $\sin(\pi x)$ ressemble à $\pm\pi(x-k)$, donc $\sin^m(\pi x)$ ressemble à $\pm\pi^m(x-k)^m$, qui s'annule à l'ordre $m$.

\begin{lemme}\label{lem:ordre-sin}
Soit $m\geq1$, $s_m(x)=\sin^m(\pi x)$ et $k\in\Z$. Alors $s_m^{(j)}(k)=0$ pour $0\leq j\leq m-1$, et $s_m^{(m)}(k)=(-1)^{km}\,m!\,\pi^m$. Autrement dit, $\ord_k(s_m)=m$.
\end{lemme}

\begin{proof}
Posons $t=x-k$. Comme $\sin(\pi x)=\sin(\pi t+\pi k)=(-1)^k\sin(\pi t)$, on a $s_m(x)=(-1)^{km}\sin^m(\pi t)$. La fonction $\sigma(t)=\frac{\sin(\pi t)}{\pi t}$, prolongée par $\sigma(0)=1$, est de classe $\Cc^\infty$ sur $\R$ : c'est la somme de la série entière $\sum_{n\geq0}\frac{(-1)^n(\pi t)^{2n}}{(2n+1)!}$, de rayon infini. Donc
\[
  s_m(x)=(-1)^{km}\pi^m\,t^m\psi(t), \qquad \psi=\sigma^m\in\Cc^\infty(\R),\quad \psi(0)=1 .
\]
Par la formule de Leibniz, pour $j\in\N$,
\[
  \frac{d^j}{dt^j}\big(t^m\psi(t)\big)\Big|_{t=0}=\sum_{i=0}^{j}\binom{j}{i}\,\big(t^m\big)^{(i)}(0)\,\psi^{(j-i)}(0).
\]
Or $\big(t^m\big)^{(i)}(0)$ vaut $m!$ si $i=m$ et $0$ sinon. La somme est donc nulle pour $j<m$, et vaut $m!\,\psi(0)=m!$ pour $j=m$. Comme $\frac{d}{dx}=\frac{d}{dt}$, on obtient le résultat.
\end{proof}

On peut maintenant répondre à la troisième question.

\begin{corollaire}\label{cor:fm}
Pour $m\geq1$, la fonction $f_m\colon x\mapsto E(x)\sin^m(\pi x)$ est de classe $\Cc^{m-1}$ sur $\R$ et de classe $\Cc^\infty$ sur $\R\setminus\Z$. Elle n'est $m$ fois dérivable en aucun entier : en $k\in\Z$, $f_m^{(m-1)}$ admet les dérivées
\[
  \big(f_m^{(m-1)}\big)'_g(k)=(-1)^{km}(k-1)\,m!\,\pi^m, \qquad \big(f_m^{(m-1)}\big)'_d(k)=(-1)^{km}k\,m!\,\pi^m .
\]
\end{corollaire}

\begin{proof}
Soit $k\in\Z$ et $I_k=]k-1,k+1[$. Sur $I_k$, $f_m$ vaut $(k-1)s_m$ à gauche de $k$ et $ks_m$ à partir de $k$. On applique le théorème \ref{thm:recollement} avec $p=(k-1)s_m$ et $q=ks_m$. La différence $p-q=-s_m$ est d'ordre $m$ en $k$ par le lemme \ref{lem:ordre-sin}. Donc $f_m$ est de classe $\Cc^{m-1}$ sur $I_k$, et $f_m^{(m-1)}$ a pour dérivées à gauche et à droite en $k$ les nombres $p^{(m)}(k)=(k-1)s_m^{(m)}(k)$ et $q^{(m)}(k)=ks_m^{(m)}(k)$. Les intervalles $I_k$ recouvrent $\R$, et être de classe $\Cc^{m-1}$ est une propriété locale, donc $f_m\in\Cc^{m-1}(\R)$. Hors de $\Z$, $f_m$ coïncide localement avec un multiple de $s_m$.
\end{proof}

\begin{remarque}
Pour $m=1$, on retrouve le corollaire \ref{cor:Esin}. Pour $m=2$, la fonction $E(x)\sin^2(\pi x)$ est de classe $\Cc^1$ mais pas deux fois dérivable aux entiers ; l'exercice 2 le vérifie à la main. La famille $(f_m)$ fournit donc, pour tout $n\in\N$, une fonction explicite de classe $\Cc^n$ et pas de classe $\Cc^{n+1}$ : il suffit de prendre $f_{n+1}$. Un exemple plus simple est $x\mapsto x^{n}|x|$, de classe $\Cc^n$ sans être $n+1$ fois dérivable en $0$. C'est le recollement de $p=-x^{n+1}$ et $q=x^{n+1}$, dont la différence $-2x^{n+1}$ est d'ordre $n+1$ en $0$.
\end{remarque}

% ==========================================================
\section{Pourquoi une dérivée ne saute jamais : le théorème de Darboux}
% ==========================================================

Regardons $f_2$ de plus près. Elle est dérivable partout, et sa dérivée $f_2'$ est continue : le saut de $E$ ne passe pas à $f_2'$. Il réapparaît un cran plus haut, dans $f_2''$, et c'est la dérivabilité de $f_2'$ qui disparaît. On pourrait espérer une situation intermédiaire : une fonction dérivable dont la dérivée saute. Cette section montre que c'est impossible.

L'outil est le théorème des accroissements finis. Il remplace un taux d'accroissement par une valeur de la dérivée en un point intermédiaire. Si la dérivée a une limite, le taux d'accroissement a la même.

\begin{proposition}[Limite de la dérivée]\label{prop:limite-derivee}
Soit $\varphi$ continue sur $[a,b[$, dérivable sur $]a,b[$, telle que $\varphi'(x)\to\ell\in\R$ quand $x\to a^+$. Alors $\varphi$ est dérivable à droite en $a$ et $\varphi'_d(a)=\ell$. On a le même énoncé à gauche.
\end{proposition}

\begin{proof}
Pour $x\in]a,b[$, le théorème des accroissements finis sur $[a,x]$ fournit $c_x\in]a,x[$ tel que $\frac{\varphi(x)-\varphi(a)}{x-a}=\varphi'(c_x)$. Quand $x\to a^+$, on a $c_x\to a^+$, donc ce taux tend vers $\ell$.
\end{proof}

\begin{corollaire}\label{cor:pas-de-saut}
Soit $\varphi$ dérivable sur un intervalle ouvert $I$ et $a\in I$. Si $\varphi'$ admet des limites finies $\varphi'(a^-)$ et $\varphi'(a^+)$, alors $\varphi'(a^-)=\varphi'(a^+)=\varphi'(a)$. Une dérivée n'a donc pas de discontinuité de première espèce.
\end{corollaire}

\begin{proof}
Par la proposition \ref{prop:limite-derivee}, $\varphi'_g(a)=\varphi'(a^-)$ et $\varphi'_d(a)=\varphi'(a^+)$. Ces deux nombres valent $\varphi'(a)$, puisque $\varphi$ est dérivable en $a$.
\end{proof}

\begin{remarque}
Ce corollaire redonne la non-dérivabilité du corollaire \ref{cor:fm}, sans calculer de taux d'accroissement. Sur $]k-1,k[$ et sur $]k,k+1[$, $f_m^{(m)}$ vaut $(k-1)s_m^{(m)}$ et $ks_m^{(m)}$. Ces deux fonctions ont en $k$ les limites $(k-1)s_m^{(m)}(k)$ et $ks_m^{(m)}(k)$, qui sont distinctes. Si $f_m^{(m-1)}$ était dérivable en $k$, sa dérivée aurait un saut en $k$, ce qui est impossible.
\end{remarque}

Le corollaire \ref{cor:pas-de-saut} suppose que les limites de $\varphi'$ existent. Darboux montre qu'on peut s'en passer : une dérivée vérifie toujours le théorème des valeurs intermédiaires, comme une fonction continue, même quand elle ne l'est pas. L'idée de la preuve : pour atteindre la valeur $\lambda$, on cherche un point où la fonction $\varphi(x)-\lambda x$ a une tangente horizontale, c'est-à-dire un minimum intérieur.

\begin{theoreme}[Darboux]\label{thm:darboux}
Soit $\varphi$ dérivable sur un intervalle $I$. Alors $\varphi'(I)$ est un intervalle : $\varphi'$ vérifie la propriété des valeurs intermédiaires.
\end{theoreme}

\begin{proof}
Soit $a<b$ dans $I$ et $\lambda$ strictement compris entre $\varphi'(a)$ et $\varphi'(b)$. Quitte à remplacer $\varphi$ par $-\varphi$ et $\lambda$ par $-\lambda$, on suppose $\varphi'(a)<\lambda<\varphi'(b)$. La fonction $\psi(x)=\varphi(x)-\lambda x$ est continue sur $[a,b]$, donc elle y atteint son minimum en un point $c$.
\begin{itemize}[nosep]
\item Comme $\psi'(a)=\varphi'(a)-\lambda<0$, on a $\psi(x)<\psi(a)$ pour $x>a$ assez proche de $a$ ; donc $c\neq a$.
\item Comme $\psi'(b)>0$, on a $\psi(x)<\psi(b)$ pour $x<b$ assez proche de $b$ ; donc $c\neq b$.
\end{itemize}
Ainsi $c\in]a,b[$ est un minimum intérieur, donc $\psi'(c)=0$, c'est-à-dire $\varphi'(c)=\lambda$.
\end{proof}

Une dérivée peut quand même être discontinue, mais pas en sautant : il faut qu'elle oscille. Voici l'exemple classique.

\begin{exemple}\label{ex:x2sin}
Soit $\varphi(x)=x^2\sin(1/x)$ pour $x\neq0$ et $\varphi(0)=0$. Comme $\big|\frac{\varphi(x)}{x}\big|\leq|x|$, $\varphi$ est dérivable en $0$, avec $\varphi'(0)=0$. Pour $x\neq0$, $\varphi'(x)=2x\sin(1/x)-\cos(1/x)$. Le premier terme tend vers $0$. Le second n'a pas de limite en $0$ : il vaut $-1$ en $x=\frac1{2n\pi}$ et $1$ en $x=\frac1{(2n+1)\pi}$. Donc $\varphi'$ est discontinue en $0$, mais d'une discontinuité de seconde espèce, comme l'imposent le corollaire \ref{cor:pas-de-saut} et le théorème de Darboux. Les fonctions $f_m$ illustrent l'autre mécanisme : la dérivée d'ordre $m$ a un vrai saut, et c'est la dérivabilité d'ordre $m$ qui disparaît.
\end{exemple}

% ==========================================================
\section{Retour à la question de départ}
% ==========================================================

On peut maintenant relire l'exercice de terminale.
\begin{itemize}[nosep]
\item \emph{Continuité.} $E$ saute aux entiers, mais $\sin(\pi x)$ s'y annule. Une fonction bornée multipliée par une fonction qui tend vers $0$ tend vers $0$ (proposition \ref{prop:produit}).
\item \emph{Points anguleux.} $\sin(\pi x)$ s'annule avec la pente $\pm\pi$. L'écart entre les pentes à droite et à gauche est le saut de $E$ multiplié par cette pente : en $0$, $1\times\pi$ ; en $1$, $1\times(-\pi)$.
\item \emph{Plus de régularité.} Chaque ordre d'annulation du facteur fait gagner un cran : $E(x)\sin^m(\pi x)$ est de classe $\Cc^{m-1}$ (corollaire \ref{cor:fm}).
\item \emph{Pas de cran intermédiaire.} On ne peut pas obtenir « dérivable avec une dérivée qui saute », car une dérivée ne saute jamais (corollaire \ref{cor:pas-de-saut} et théorème de Darboux).
\end{itemize}

% ==========================================================
\section{Exercices}
% ==========================================================

\begin{exercice}
Déterminer les points de continuité et de dérivabilité sur $\R$ de $g\colon x\mapsto E(x)\big(x-E(x)\big)$.
\end{exercice}

\begin{exercice}
Soit $f_2(x)=E(x)\sin^2(\pi x)$. Montrer, sans utiliser le théorème \ref{thm:recollement}, que $f_2$ est dérivable sur $\R$ et que $f_2'(x)=\pi E(x)\sin(2\pi x)$ pour tout réel $x$. En déduire que $f_2$ est de classe $\Cc^1$, et calculer les dérivées à gauche et à droite de $f_2'$ en $k\in\Z$.
\end{exercice}

\begin{exercice}
Déterminer les réels $\alpha,\beta,\gamma$ pour lesquels la fonction $\varphi$, égale à $e^x$ sur $]-\infty,0[$ et à $\alpha+\beta x+\gamma x^2$ sur $[0,+\infty[$, est de classe $\Cc^2$ sur $\R$. Est-elle alors trois fois dérivable ?
\end{exercice}

\begin{exercice}
Montrer qu'il n'existe aucune fonction $F$ dérivable sur $\R$ telle que $F'=E$.
\end{exercice}

\begin{exercice}
Soit $G(x)=\int_0^xE(t)\,dt$ pour $x\in\R$.
\begin{enumerate}[nosep]
\item Montrer que $G(x)=nx-\frac{n(n+1)}2$, où $n=E(x)$.
\item Montrer que $G(x)=\max_{n\in\Z}\big(nx-\frac{n(n+1)}2\big)$, et en déduire que $G$ est convexe.
\item Déterminer les points où $G$ est dérivable, et ses dérivées à gauche et à droite aux autres points.
\end{enumerate}
\end{exercice}

% ==========================================================
\section{Corrigés}
% ==========================================================

\paragraph{Exercice 1.} \emph{Idée : ici le facteur $x-E(x)$ ne s'annule pas à gauche des sauts, donc la proposition \ref{prop:produit} ne s'applique pas, et l'on s'attend à des discontinuités.}

Sur $[k,k+1[$, $g(x)=k(x-k)$, donc $g$ est de classe $\Cc^\infty$ sur $\R\setminus\Z$. En $k\in\Z$, on a $g(k)=0$ et $g(k^-)=(k-1)\big(k-(k-1)\big)=k-1$. Donc $g$ est continue en $k$ si et seulement si $k=1$. En $1$, $g$ est nulle sur $[0,1[$ et vaut $x-1$ sur $[1,2[$, donc $g'_g(1)=0\neq1=g'_d(1)$. Ainsi $g$ est continue exactement sur $\R\setminus(\Z\setminus\{1\})$, et dérivable exactement sur $\R\setminus\Z$. L'échec de la continuité vient de ce que $x-E(x)\to1$ à gauche des entiers : seul le saut en $1$ est compensé, parce que là c'est le facteur $E(x)$ lui-même qui vaut $0$ à gauche.

\paragraph{Exercice 2.} \emph{Idée : appliquer deux fois la proposition \ref{prop:produit}, une fois à $f_2$ avec $h=\sin^2(\pi\,\cdot)$, de pente nulle aux entiers, puis à $f_2'$ avec $h=\sin(2\pi\,\cdot)$, de pente non nulle.}

Hors de $\Z$, $f_2$ coïncide localement avec $k\sin^2(\pi x)$, de dérivée $2k\pi\sin(\pi x)\cos(\pi x)=k\pi\sin(2\pi x)$. En $k\in\Z$, on applique la proposition \ref{prop:produit} avec $g=E$ et $h=\sin^2(\pi\,\cdot)$. On a $h(k)=0$ et $h'(k)=\pi\sin(2\pi k)=0$, donc $f_2$ est dérivable en $k$, de dérivée nulle, égale à $\pi E(k)\sin(2\pi k)$. Ainsi $f_2'=\pi E\cdot\sin(2\pi\,\cdot)$ sur $\R$. Cette fonction est continue par la proposition \ref{prop:produit}(1), puisque $\sin(2\pi k)=0$. Avec $h=\sin(2\pi\,\cdot)$ et $h'(k)=2\pi$, la proposition \ref{prop:produit}(2) donne
\[
  (f_2')'_g(k)=2\pi^2(k-1), \qquad (f_2')'_d(k)=2\pi^2k .
\]
C'est le corollaire \ref{cor:fm} pour $m=2$ : $(-1)^{2k}\,2!\,\pi^2=2\pi^2$. Ainsi $f_2$ est de classe $\Cc^1$, mais pas deux fois dérivable aux entiers.

\paragraph{Exercice 3.} \emph{Idée : d'après le théorème \ref{thm:recollement}, il suffit de raccorder les valeurs et les deux premières dérivées en $0$, c'est-à-dire de prendre le développement de Taylor de $e^x$ à l'ordre $2$.}

On applique le théorème \ref{thm:recollement} avec $p=\exp$, $q(x)=\alpha+\beta x+\gamma x^2$ et $a=0$. On a $(p-q)(0)=1-\alpha$, $(p-q)'(0)=1-\beta$ et $(p-q)''(0)=1-2\gamma$. Donc $\varphi\in\Cc^2(\R)$ si et seulement si $\alpha=\beta=1$ et $\gamma=\frac12$ : $q$ est alors le polynôme de Taylor d'ordre $2$ de l'exponentielle en $0$. Dans ce cas, $(p-q)'''(0)=1\neq0$, donc $\ord_0(p-q)=3$, et $\varphi$ n'est pas trois fois dérivable en $0$ : $\varphi''$ a pour dérivée à gauche $p'''(0)=1$ et pour dérivée à droite $q'''(0)=0$.

\paragraph{Exercice 4.} \emph{Idée : $E$ saute, et une dérivée ne saute pas.}

Si $F'=E$ sur $\R$, le théorème de Darboux \ref{thm:darboux} impose que $E([0,1])$ soit un intervalle. Or $E([0,1])=\{0,1\}$, qui ne contient pas $\frac12$. Plus directement, $E$ a un saut en $0$, ce que le corollaire \ref{cor:pas-de-saut} interdit. En revanche, $E$ admet des primitives en un sens plus faible, comme la fonction $G$ de l'exercice suivant, qui est continue et dérivable hors de $\Z$, de dérivée $E$.

\paragraph{Exercice 5.} \emph{Idée : $G$ est affine sur chaque $[n,n+1]$, de pente $n$ ; ses morceaux sont des droites dont la pente augmente, ce qui est exactement la convexité.}

1. Pour $n\in\Z$ et $x\in[n,n+1]$, on a $E=n$ sur $[n,x[$, donc $G(x)-G(n)=n(x-n)$. En particulier, $G(n+1)-G(n)=n$. La suite $u_n=\frac{n(n-1)}2$ vérifie $u_0=0=G(0)$ et $u_{n+1}-u_n=n$. Par récurrence vers le haut et vers le bas, $G(n)=\frac{n(n-1)}2$ pour tout $n\in\Z$. Pour $x\in[n,n+1[$, on obtient
\[
  G(x)=\frac{n(n-1)}2+n(x-n)=nx-\frac{n(n+1)}2 .
\]

2. Notons $\ell_n(x)=nx-\frac{n(n+1)}2$. On a
\[
  \ell_n(x)-\ell_{n-1}(x)=x-\frac{n(n+1)-(n-1)n}2=x-n .
\]
Soit $x\in\R$ et $N=E(x)$. Pour $n\leq N$, $\ell_n(x)-\ell_{n-1}(x)=x-n\geq0$. Pour $n\geq N+1$, $\ell_n(x)-\ell_{n-1}(x)=x-n<0$. La suite $(\ell_n(x))_{n\in\Z}$ croît donc jusqu'au rang $N$, puis décroît, et son maximum est $\ell_N(x)=G(x)$. Comme borne supérieure de fonctions affines, $G$ est convexe.

3. Sur $]n,n+1[$, $G$ est affine de pente $n$. En $n\in\Z$, $G$ coïncide avec $\ell_{n-1}$ sur $[n-1,n]$ et avec $\ell_n$ sur $[n,n+1]$, donc $G'_g(n)=n-1$ et $G'_d(n)=n$. Ainsi $G$ est dérivable exactement sur $\R\setminus\Z$. On retrouve deux propriétés générales des fonctions convexes : leur dérivée à droite est croissante (ici $G'_d=E$), et elles sont dérivables hors d'un ensemble au plus dénombrable.

\begin{thebibliography}{9}
\bibitem{gourdon} X. Gourdon, \emph{Les maths en tête. Analyse}, Ellipses.
\bibitem{hauchecorne} B. Hauchecorne, \emph{Les contre-exemples en mathématiques}, Ellipses.
\bibitem{rombaldi} J.-E. Rombaldi, \emph{Éléments d'analyse réelle}, EDP Sciences.
\end{thebibliography}

\paragraph{Source.} Partie « Continuité » de la feuille de révision de terminale S (2003), d'après les manuels \emph{Hyperbole} et \emph{Bréal} ; le reste de la feuille est traité dans la leçon 142.

\end{document}
